\selectlanguage{greek}
\begin{abstract}

Τα μοντέλα διάχυσης έχουν αναδειχθεί ως το επικρατέστερο παράδειγμα για τη
σύνθεση εικόνων υψηλής πιστότητας, ωστόσο η ανάπτυξή τους σε περιβάλλοντα
περιορισμένων υπολογιστικών πόρων παραμένει δυσχερής λόγω του υψηλού
υπολογιστικού κόστους των αρχιτεκτονικών μεγάλης χωρητικότητας. Η παρούσα
διπλωματική εργασία διερευνά την εφαρμογή της απόσταξης γνώσης σε μοντέλα
διάχυσης χώρου εικονοστοιχείων ως μηχανισμό βελτίωσης της ποιότητας παραγωγής
ενός συμπαγούς μοντέλου-μαθητή, χωρίς τροποποίηση της αρχιτεκτονικής του κατά
την εξαγωγή συμπερασμάτων ή του υπολογιστικού του αποτυπώματος --- έναν στόχο
ορθογώνιο προς το ευρέως μελετημένο πρόβλημα της επιτάχυνσης δειγματοληψίας.
Στο πλαίσιο ενός σχήματος δασκάλου--μαθητή με παγωμένες παραμέτρους δασκάλου,
βασισμένου σε αρχιτεκτονική Vision Transformer εκπαιδευμένη με αντικειμενική
συνάρτηση αντιστοίχισης ροής, προτείνουμε και αξιολογούμε συστηματικά ένα ενιαίο
σχήμα απόσταξης που συνδυάζει ευθυγράμμιση ενδιάμεσων χαρτών χαρακτηριστικών,
μεταφορά προτύπων προσοχής, στάθμιση απωλειών με επίγνωση του βήματος χρόνου,
και αρχικοποίηση βαρών από τον δάσκαλο. Πειράματα σε σημεία αναφοράς παραγωγής
εικόνων καταδεικνύουν ότι το προτεινόμενο πλαίσιο μπορεί σε κάποιες περιπτώσεις να βελτιώσει την ποιότητα παραγωγής του μοντέλου-μαθητή σε σχέση με ένα αντίστοιχο μοντέλο εκπαιδευμένο από την αρχή υπό πανομοιότυπο υπολογιστικό προϋπολογισμό, επικυρώνοντας την απόσταξη γνώσης ως αποτελεσματική στρατηγική βελτίωσης ποιότητας για συμπαγή
μοντέλα διάχυσης.

   \begin{keywords}
      Μοντέλα Διάχυσης, Απόσταξη Γνώσης, Αντιστοίχιση Ροής,
      Παραγωγή Εικόνων, Απόσταξη Χαρακτηριστικών, Μεταφορά Προσοχής
   \end{keywords}
\end{abstract}


\selectlanguage{english}

\selectlanguage{english}
\begin{abstracteng}

Diffusion models have emerged as the dominant paradigm for high-fidelity image
synthesis, yet their deployment in resource-constrained settings remains impeded by
the substantial computational cost of large-capacity architectures. This thesis
investigates the application of knowledge distillation to pixel-space diffusion models
as a mechanism for improving the generative quality of a compact student model,
without modifying its inference-time architecture or computational footprint --- an
objective orthogonal to the well-studied problem of sampling acceleration. Operating
within a frozen teacher--student framework built upon a Vision Transformer backbone
trained with a flow-matching objective, we propose and systematically evaluate a
unified distillation scheme that combines intermediate feature map alignment,
attention pattern transfer, timestep-aware loss weighting, and teacher-derived weight
initialisation. Experiments on image generation benchmarks demonstrate that the
proposed framework under specific configurations can improve the generative quality of the student model
relative to a matched baseline trained from scratch under an identical computational
budget, validating knowledge distillation as an effective quality-boosting strategy for
compact diffusion models.

   \begin{keywordseng}
      Diffusion Models, Knowledge Distillation, Flow Matching, Vision Transformers,
      Image Generation, Feature Distillation, Attention Transfer
   \end{keywordseng}
\end{abstracteng}
